% Part: first-order-logic
% Chapter: axiomatic-deduction
% Section: provability-quantifiers

% verification of properties of provability needed for maximally
% consistent sets in the completeness chapter.

\documentclass[../../../include/open-logic-section]{subfiles}

\begin{document}

\olfileid[hi]{fol}{axd}{qpr}

\olsection{\usetoken{S}{derivability} और परिमाणक}

\begin{explain}
  पूर्णता प्रमेय के लिए यह भी आवश्यक है कि अभिगृहीतीय व्युत्पत्तियाँ इस
  अनुभाग में स्थापित~$\Proves$ से संबंधित तथ्य प्रदान करें।
\end{explain}

\begin{thm}
\ollabel{thm:strong-generalization} यदि $c$ ऐसा !!a{constant} है जो
$\Gamma$ या $!A(x)$ में नहीं आता, और $\Gamma \Proves !A(c)$, तो $\Gamma
\Proves \lforall[x][!A(x)]$।
\end{thm}

\begin{proof}
निगमन प्रमेय से $\Gamma \Proves \ltrue \lif !A(c)$। चूँकि $c$, $\Gamma$
या~$\top$ में नहीं आता, इसलिए $\Gamma \Proves
\ltrue \lif \lforall[x][!A(x)]$ मिलता है। निगमन प्रमेय से फिर $\Gamma \Proves
\lforall[x][!A(x)]$।
\end{proof}

\begin{prop}
\ollabel{prop:provability-quantifiers}
\begin{tagenumerate}{prvEx,prvAll}
\tagitem{prvEx}{$!A(t) \Proves \lexists[x][!A(x)]$.}{}

\tagitem{prvAll}{$\lforall[x][!A(x)] \Proves !A(t)$.}{}
\end{tagenumerate}
\end{prop}

\begin{proof}
\begin{tagenumerate}{prvEx,prvAll}
\tagitem{prvEx}{\olref[qua]{ax:q2} और निगमन प्रमेय से।}{}
\tagitem{prvAll}{\olref[qua]{ax:q1} और निगमन प्रमेय से।}{}
\end{tagenumerate}
\end{proof}

\end{document}
